\section{Experiments}
% ======================================================================

\problem{experiment\_log}{Experiment logging (3 points)}

\deliverable{Logging infrastructure code for your experiments and an experiment log for the problems below.}

\textit{Answer:} % TODO: describe your logging infrastructure and experiment log


\problem{learning\_rate}{Tune the learning rate (3 points)}

\textbf{(a)} Perform a hyperparameter sweep over learning rates and report the final losses.

\deliverable{Learning curves associated with multiple learning rates. Explain your hyperparameter search strategy.}

\begin{figure}[H]
    \centering
    % \includegraphics[width=0.8\textwidth]{figures/learning_rate_sweep.png}
    \fbox{\parbox{0.8\textwidth}{\centering\vspace{2cm} Learning rate sweep learning curves \\ \textit{(insert figure here)}\vspace{2cm}}}
    \caption{Learning curves for different learning rates on TinyStories.}
    \label{fig:lr_sweep}
\end{figure}

\textit{Search strategy:} % TODO: describe your strategy

\deliverable{A model with validation loss (per-token) on TinyStories of at most 1.45.}

\textit{Achieved validation loss:} % TODO: report your loss

\textbf{(b)} Investigate how the point at which learning rates diverge is related to your best learning rate.

\deliverable{Learning curves of increasing learning rate with at least one divergent run, and an analysis.}

\begin{figure}[H]
    \centering
    % \includegraphics[width=0.8\textwidth]{figures/learning_rate_divergence.png}
    \fbox{\parbox{0.8\textwidth}{\centering\vspace{2cm} Learning rate divergence analysis \\ \textit{(insert figure here)}\vspace{2cm}}}
    \caption{Learning curves showing divergence at high learning rates.}
    \label{fig:lr_divergence}
\end{figure}

\textit{Analysis:} % TODO: your analysis


\problem{batch\_size\_experiment}{Batch size variations (1 point)}

\deliverable{Learning curves for runs with different batch sizes.}

\begin{figure}[H]
    \centering
    % \includegraphics[width=0.8\textwidth]{figures/batch_size_experiment.png}
    \fbox{\parbox{0.8\textwidth}{\centering\vspace{2cm} Batch size experiment learning curves \\ \textit{(insert figure here)}\vspace{2cm}}}
    \caption{Learning curves for different batch sizes on TinyStories.}
    \label{fig:batch_size}
\end{figure}

\deliverable{A few sentences discussing your findings on batch sizes and their impacts on training.}

\textit{Findings:} % TODO: your findings


\problem{generate}{Generate text (1 point)}

\deliverable{Text dump of at least 256 tokens of text (or until the first \texttt{<eos>} token), and a brief comment on fluency.}

\textit{Generated text:}

\begin{quote}
% TODO: paste your generated text here (at least 256 tokens)
\end{quote}

\textit{Fluency commentary:} % TODO: comment on fluency and at least two factors


\problem{layer\_norm\_ablation}{Remove RMSNorm and train (1 point)}

\deliverable{A learning curve for when you remove RMSNorms and train, as well as a learning curve for the best learning rate.}

\begin{figure}[H]
    \centering
    % \includegraphics[width=0.8\textwidth]{figures/layer_norm_ablation.png}
    \fbox{\parbox{0.8\textwidth}{\centering\vspace{2cm} RMSNorm ablation learning curves \\ \textit{(insert figure here)}\vspace{2cm}}}
    \caption{Learning curves comparing with and without RMSNorm.}
    \label{fig:layer_norm_ablation}
\end{figure}

\deliverable{A few sentences of commentary on the impact of RMSNorm.}

\textit{Commentary:} % TODO: your commentary


\problem{pre\_norm\_ablation}{Implement post-norm and train (1 point)}

\deliverable{A learning curve for a post-norm Transformer, compared to the pre-norm one.}

\begin{figure}[H]
    \centering
    % \includegraphics[width=0.8\textwidth]{figures/pre_norm_ablation.png}
    \fbox{\parbox{0.8\textwidth}{\centering\vspace{2cm} Pre-norm vs.\ post-norm learning curves \\ \textit{(insert figure here)}\vspace{2cm}}}
    \caption{Learning curves comparing pre-norm and post-norm Transformers.}
    \label{fig:pre_norm_ablation}
\end{figure}


\problem{no\_pos\_emb}{Implement NoPE (1 point)}

\deliverable{A learning curve comparing the performance of RoPE and NoPE.}

\begin{figure}[H]
    \centering
    % \includegraphics[width=0.8\textwidth]{figures/nope_vs_rope.png}
    \fbox{\parbox{0.8\textwidth}{\centering\vspace{2cm} RoPE vs.\ NoPE learning curves \\ \textit{(insert figure here)}\vspace{2cm}}}
    \caption{Learning curves comparing RoPE and NoPE.}
    \label{fig:nope_vs_rope}
\end{figure}


\problem{swiglu\_ablation}{SwiGLU vs.\ SiLU (1 point)}

\deliverable{A learning curve comparing the performance of SwiGLU and SiLU feed-forward networks.}

\begin{figure}[H]
    \centering
    % \includegraphics[width=0.8\textwidth]{figures/swiglu_vs_silu.png}
    \fbox{\parbox{0.8\textwidth}{\centering\vspace{2cm} SwiGLU vs.\ SiLU learning curves \\ \textit{(insert figure here)}\vspace{2cm}}}
    \caption{Learning curves comparing SwiGLU and SiLU feed-forward networks.}
    \label{fig:swiglu_vs_silu}
\end{figure}

\deliverable{A few sentences discussing your findings.}

\textit{Findings:} % TODO: your findings


\problem{main\_experiment}{Experiment on OWT (2 points)}

\deliverable{A learning curve of your language model on OpenWebText. Describe the difference in losses from TinyStories.}

\begin{figure}[H]
    \centering
    % \includegraphics[width=0.8\textwidth]{figures/owt_training.png}
    \fbox{\parbox{0.8\textwidth}{\centering\vspace{2cm} OpenWebText training learning curve \\ \textit{(insert figure here)}\vspace{2cm}}}
    \caption{Learning curve for the language model trained on OpenWebText.}
    \label{fig:owt_training}
\end{figure}

\textit{Loss comparison:} % TODO: describe the difference in losses

\deliverable{Generated text from OpenWebText LM. How is the fluency? Why is the output quality worse?}

\textit{Generated text:}

\begin{quote}
% TODO: paste your generated text here
\end{quote}

\textit{Fluency commentary:} % TODO: comment on fluency and why it's worse


\problem{leaderboard}{Leaderboard (6 points)}

\deliverable{The final validation loss, an associated learning curve (wall-clock-time x-axis, < 45 min), and a description of what you did.}

\begin{figure}[H]
    \centering
    % \includegraphics[width=0.8\textwidth]{figures/leaderboard.png}
    \fbox{\parbox{0.8\textwidth}{\centering\vspace{2cm} Leaderboard submission learning curve \\ \textit{(insert figure here)}\vspace{2cm}}}
    \caption{Learning curve for the leaderboard submission (wall-clock time on x-axis).}
    \label{fig:leaderboard}
\end{figure}

\textit{Final validation loss:} % TODO: report your loss

\textit{Description:} % TODO: describe what modifications you made

